% Generated draft: requires mathematical and source review.
\documentclass[11pt]{article}
\usepackage[margin=25mm]{geometry}
\usepackage{fontspec}
\setmainfont{Latin Modern Roman}
\usepackage{amsmath,amssymb,mathtools,unicode-math}
\setmathfont{Latin Modern Math}
\usepackage{enumitem,needspace,xurl,hyperref}
\setlength{\parindent}{0pt}
\setlength{\parskip}{5pt}
\setlength{\emergencystretch}{3em}
\newcommand{\N}{\mathbb{N}}
\newcommand{\Z}{\mathbb{Z}}
\newcommand{\Q}{\mathbb{Q}}
\newcommand{\R}{\mathbb{R}}
\newcommand{\C}{\mathbb{C}}
\newcommand{\F}{\mathbb{F}}
\newcommand{\A}{\mathbb{A}}
\newcommand{\PP}{\mathbb{P}}
\newcommand{\E}{\mathbb{E}}
\newcommand{\eps}{\varepsilon}
\newcommand{\dd}{\,\mathrm{d}}
\begin{document}
\section*{UnsolvedMath ID: 20003343}
\textbf{Source problem code:} catalogue entry
\subsection{Definitions and mathematical statement}
For an arithmetic homogeneous space associated with a semisimple $\Q$-group $\mathbf G$, let $r_{\Q}=\operatorname{rank}_{\Q}\mathbf G$ and let $D$ be the specified connected real split torus. In the usual full-orbit interpretation, a $D$-orbit is divergent if its orbit map $D\to G/\Gamma$ is proper: every sequence escaping compact subsets of $D$ has images escaping compact subsets of the quotient. The source's preceding definitions are absent, including the precise torus, lattice and meaning of ``obvious'' divergence.

B. Weiss's Conjecture 48 prints three claims: if $\dim D>r_{\Q}$, there are no divergent orbits; if $\dim D=r_{\Q}$, only obvious divergence occurs; if $\dim D<r_{\Q}$, non-obvious divergent trajectories exist. The aid reports clause (1) proved by Weiss, clause (2) by Solan--Tamam, and clause (3) by Tamam for almost $\Q$-simple groups, with a factorwise correction needed in general. The exact theorems and corrected criterion are not supplied. These claims must be separated: literal clause (3) is reported false for unrestricted semisimple products.

The aid's counterexample family fixes any real quadratic field $K$ and integer $n\geq3$. Put
\[
 \mathbf H=\operatorname{Res}_{K/\Q}SL_2,\qquad
 \mathbf G=\mathbf H\times SL_n,
 \qquad G=\mathbf G(\R),
\]
with the product arithmetic lattice
\[
 \Gamma=SL_2(\mathcal O_K)\times SL_n(\Z),
\]
where the first factor is embedded via both real embeddings of $K$. The real group $\mathbf H(\R)$ is $SL(2,\R)\times SL(2,\R)$, with real split rank two but $\Q$-rank one. Let $D_H$ be its two-dimensional maximal connected real split torus and let $D=D_H\times\{e\}$ act trivially on the $SL_n$ factor. Then
\[
 \dim D=2<\operatorname{rank}_{\Q}\mathbf G
 =1+(n-1)=n.
\]
Thus the global inequality required by printed clause (3) holds.

For product lattices the quotient is a product, and the orbit map has form $d\mapsto(dx_1,x_2)$. The reported product-properness lemma says it is proper exactly when $d\mapsto dx_1$ is proper in the active factor. In that factor,
\[
 \dim D_H=2>\operatorname{rank}_{\Q}\mathbf H=1,
\]
so the cited Weiss theorem excludes every divergent $D_H$ orbit. Hence the product has no divergent $D$ orbit at all, and in particular no non-obvious one, despite $\dim D<\operatorname{rank}_{\Q}\mathbf G$.

The inactive factor raises total rational rank without contributing dynamics. This supplies the reported literal-product counterexample, not a refutation of the almost-$\Q$-simple theorem or an unstated factorwise version. It does not define the source's obvious-divergence class or supply the full repaired trichotomy; those require the preceding setup and full references. The solved/refuted registry label applies to this clause-specific reading rather than asserting that all three printed clauses are true.
\subsection{Short English statement}
The three clauses have different reported outcomes. A product factor on which the torus acts trivially can increase total rational rank while the active factor still forbids divergence. The explicit family refutes the unrestricted third clause, while leaving the almost-simple and corrected factorwise statements distinct.
\subsection{Sources}
\begin{itemize}
\item Local UnsolvedMath registry, permanent ID 20003343, source code catalogue entry; \texttt{aim-other-notes.json}, notes[149], source tag conjecture.
\item AIM, \emph{Emerging applications of measure rigidity}, Conjecture 48, B. Weiss: \url{https://aimath.org/WWN/measrigid/measrigid.pdf}. Definitions from the preceding record are not supplied here.
\item Weiss, Solan--Tamam and Tamam are credited by the research summary for the respective clause results; no full bibliography or corrected factorwise theorem statement is supplied.
\item Machine-generated research attempt 1: counterexample, refuted at the run, rigor labelled ``proved,'' low novelty confidence; registry solved/``SOLVED-IN-LITERATURE.'' Not peer reviewed; independently verify before citation. Registry updated 14 August 2026.
\end{itemize}
\end{document}
